\section{The genocide exception}
\label{sec:F.4}
The hardest case for this term is force to stop a genocide when the Council is blocked by a veto, and it runs on the extended term. Rwanda in 1994 and Kosovo in 1999 are the references. The prohibition on genocide is peremptory, and the Genocide Convention imposes a duty to prevent \autocite{icj2007bosnia}. My reason for discounting the Council's judgment on defense claims applies here with more force: the permanent member shielding a client has the largest stake of all. “Humanitarian intervention” is also a phrase wars of choice are fought under, so the term can't lift on the intervener's claim, and it needn't: the device that governs the defense exception governs this one. The exception is the slow route by design. It licenses a machine \emph{into} a pipeline for that purpose, and waits on findings. A machine already present that can stop an atrocity under way acts under the one permission, and waits on nothing but the scene. So the exception's slowness costs less than it looks, and it can afford the rule that follows. The term lifts for force to stop a genocide under way on a merits judgment of the International Court of Justice under the Genocide Convention, or on two findings of different kinds, at least one of which comes from a body that took evidence, a court or a commission. The kinds are the four in \ref{sec:F.1}, and for this exception a fifth: a provisional-measures order of the Court finding the claim plausible and the risk real. A Council or Assembly resolution is a vote. It corroborates, and two votes never make the pair. A provisional-measures order counts here and nowhere else in the term, and only as one of two: an adversarial court's finding of plausibility and an investigative body's published evidence are unlike sources converging, which is the strongest corroboration the discount and the burden know. The named state's response to a finding, or its refusal to respond, goes on the record beside the finding, and refusal reads as withheld evidence. What I condemn in a finding about a person is that it is unreasoned, unpublished and unappealable. A commission publishes reasons and evidence, invites a response, and is displaced by a merits judgment. It fails only on neutrality between the parties, and the two-body rule answers that. It lifts only for the force and the place the finding covers. The force is direct: against the forces committing the genocide, where they are committing it, and the least that stops it. Force elsewhere against the state responsible, on the theory that the cost will make it stop, is outside the exception whatever its motive. The theory is a counterfactual supplied by the intervener, the class of claim hardest to corroborate, and a licence for it would be a licence for a general war under a humanitarian heading. It can license a machine on the other side of a war whose defender holds an emergency licence for a different act, and the person-directed terms and the law of armed conflict govern both. Without this exception the term would refuse where refusing leaves the people being destroyed with no one answerable to them at all, in cases identifiable in advance, with a remedy the term already uses. A floor that protects standing can't be the reason a population is erased from every relation of standing it had.

A finding under the genocide exception also counts against the belligerent it names. Force a disinterested body has found to be genocide is not force repelling an attack. The finding lifts that belligerent's emergency licence for the force and place it covers, and leaves whatever remains that repels an attack, which may be nothing.
